% !TeX program = pdflatex

\documentclass[final]{ltugboat}

\usepackage{mwe}
\usepackage[hyphens]{url}
\usepackage[hidelinks]{hyperref}
\usepackage{xspace}

\def\JR{Jab\-Ref\xspace}
\def\SLR{\acro{SLR}}

\title{JabRef: \BibTeX-based literature management~software}

\author{Oliver Kopp}
\address{Sindelfingen, Germany}
\EDITORnonetaddress
\personalURL{https://github.com/koppor}
\ORCID{0000-0001-6962-4290}

\author{Carl Christian Snethlage}
\address{Zwiesel, Germany}
\EDITORnonetaddress
\personalURL{https://github.com/calixtus}

\author{Christoph Schwentker}
\address{Paderborn, Germany}
\EDITORnonetaddress
\personalURL{https://github.com/Siedlerchr}

\begin{document}
\maketitle

\begin{abstract}
\JR\ is a comprehensive literature management tool built entirely around the \BibTeX\ data format.
This excerpt focuses on \JR's support for systematic literature reviews and serves as a
manual test fixture for extracting related-work text from \PDF{}s.
\end{abstract}

\section{Introduction}

\blindtext

\section{Related Work}

Systematic literature reviews (\SLR{}s) are comprehensive examinations of existing research on a particular topic, conducted according to a well-defined and transparent methodology~\cite{Kitchenham2007}.
The most difficult challenge researchers face when conducting an \SLR\ is during the search step~\cite{AlZubidy2018,Voigt2021}.
Commonly used e-libraries in the domain of computer science research, such as \acro{IEEE}, arXiv, and \acro{ACM}, do not support easy bulk access, which is key to the \SLR\ method~\cite{AlZubidy2018}.
\JR fills this gap with its special \SLR\ feature~\cite{Voigt2021}.
The main concept of the implementation is to a)~build on git and \BibTeX{} to track the state of papers (excluded, included, \dots); and b)~transform a generic query syntax to the different query formats of each publisher.

\section{Conclusion}

\blindtext

\bibliographystyle{tugboat}
\bibliography{references}

\makesignature

\end{document}
